\documentclass[11pt,a4paper,twoside]{article}
\usepackage[utf8]{inputenc}
\usepackage[english]{babel}
\usepackage{actuarialangle, actuarialsymbol}
\usepackage{amsmath,amsfonts,amsthm,amssymb,mathtools}

\makeatletter
\newcommand{\ix}{\@actothersc{i}}
\makeatother

\begin{document}

Give an overview of the factors involved composed population models (pension insurance)

INSERT IMAGE FROM SLIDE 15

Additionaly we require the Axioms of Pension Insurance Mathematics (Neuburger):
1. Eliminations from populations (within one year) are uniformly distributed
2. Linear inter-annual interest with credit at the end of the year
3. Pension payments occur at beginning or end of periods


Use the following notation 
We consider the events disability and death \\
\( T_\gamma \) : Time of event \( \gamma \) \\
\( X_\gamma = T_\gamma - t_0 \): Age at event  \( \gamma \) \\
let \( \gamma = 1 \) be the beginning of disability and \( \gamma = 2 \) death. \\
We can write explicitely the additional probability terms at age \( x \):

Elimination from the population of actives within one year by becoming disabled
\[ \actsymb{i}{x} = \mathbb{P}[X_1 \leq \min(x+1, X_2) | X_1 > x, X_2 > x] \]
Elimination from the population of actives within one year by death
\[ \qx{x}[aa] = \mathbb{P}[X_2 < \min(x+1, X_1) | X_1 > x, X_2 > x ] \]
Being in the population of actives (unconditional)
\[ l^a_x = \mathbb{P}[X_1 > x, X_2 > x] \]
Being in the population of actives (conditional)
\[ \mathbb{P}[X_1 > x|X_2 > x] = \frac{\mathbb{P}[X_1 > x, X_2 > x]}{ \mathbb{P}[X_2 > x] } = \frac{l^a_x}{l_x} \]
Being in the population of disabled 
\[ \mathbb{P}[X_1 \leq x | X_2 > x] = 1 - \frac{l^a_x}{l_x} = \frac{l_x - l_x^a}{l_x} \]
Remaining in the population of actives
\[ \px{x}[a] = 1 - (\actsymb{i}{x} + \qx{x}[aa]) = \mathbb{P}[X_1 > x+1, X_2 > x+1|X_1 > x, X_2 > x] \]
Disabled population mortality
\[ \qx{x}[i] = \mathbb{P}[X_2 \leq x+1|X_1 \leq x, X_2 > x] = \qx{x}[ai] + \qx{x}[aa] \]
Active population mortality
\[ \qx{x}[a] = \mathbb{P}[X_2 \leq x+1|X_1 > x, X_2 > x] \]
Total population mortality
\[ \qx{x} = \mathbb{P}[X_2 \leq x+1|X_2 > x] \]
Total population survival
\[ \px{x} = \mathbb{P}[X_2 > x+1|X_2 > x] \]
Total population survival until \( x \)
\[ l_x = \mathbb{P}[X_2 > x] \]

Moreover we have the identities: \\
Consistency Equation.
\[ \qx{x} = \qx{x}[i] \mathbb{P}[X_1 \leq x | X_2 > x] + \qx{x}[a] \mathbb{P}[X_1 > x|X_2 > x] \]
Under the assumption that transitions between populations occur midyear (almost done exclusively in practice):
\[ \qx{x} = \qx{x}[i] - \frac{l^a_x}{l_x}(\qx{x}[i] - \qx{x}[aa] - \actsymb{i}{x} \qx[\frac{1}{2}]{x + \frac{1}{2}}[i])  \]
Keep in mind that by axiom (2) we have for \( u \in [0, 1] \) , \( \qx[u]{x}[\gamma] = u \qx[u]{x}[\gamma] \)
for example \( u = \frac{1}{2} \) then
\[( \qx[\frac{1}{2}]{x + \frac{1}{2}} = 1 - \frac{1 - q}{1 - \qx[\frac{1}{2}]{x}} = \frac{\frac{1}{2}\qx{x}}{1 - \frac{1}{2}\qx{x}} \]

Note that this analysis can be extended by changing the probabilities with respect to cohorts (different years of birth \( t_0 \)
and not assuming stationarity for mortality tables).


Describe surviving dependents in the composed population model (pension insurance)


\(h_x\) denotes the probability of a person at age \( x \) to be married at
time of death before age \( x + 1 \). \\
Let \( y(x) \) be the average age of the surviving dependent, given death of the insured
person between age \( x \) and \( x+1 \) \\
Let \( \qx{y}[w] \) be the probability of a \( y \)-year old surviving dependent to die
within one year. \\
Keep in mind that a surviving dependent might be born in a different cohort which could have a different mortality table.
We will supress this in the notation, but its straightforward to adapt the formulas.\\

Introducing the related probabilities \\
Disabled within one year and die at the end of the year as a disabled
\[ \qx{x}[ai] = \actsymb{i}{x} \qx[\frac{1}{2}]{x + \frac{1}{2}}[i] \]
Disabled within one year and survive the end of the year as a disabled
\[ \px{x}[ai] = \actsymb{i}{x} \px[\frac{1}{2}]{x + \frac{1}{2}}[i] \]
Mortality of an active
\[ \qx{x}[a] = \qx{x}[aa] = \actsymb{i}{x} \qx[\frac{1}{2}]{x + \frac{1}{2}}[i]\]

Probability of a person of age \( x \) to die within one year and to have a surviving
dependent who survives until the end of the year with age \( y_x + 1 \)
\[ \px{x}[\gamma w] = \qx{x}h_x \px[\frac{1}{2}]{y_x + \frac{1}{2}}[w] \]
where \( \gamma \) refers to the respective subpopulation. \\

It holds the second consistency equation 
\[ \px{x}[w] = \frac{l^a_x}{l_x}\px{x}[aw] + \frac{l_x - l_x^a}{l_x}\px{x}[iw] \]
with the identity \( \px{x}[aw] = \px{x}[aaw] + \px{x}[aiw] \)


Value pension insurance products and explain the two different approaches for widows

Let
\( M = [X_1] - x \) Number of complete years until the event “Disability”
\(N = [X_2] - x\): Number of complete years until the event “Death
Disability pension:\\
Settlement value:
\[ B = v^{M+1} \ax**{\angl{N - M}} \mathbf{1}_{\{N > M\}} \]
Expected value
\[ \mathbb{E}[B] = \ax**{x}[ai] = \sum_{k \geq 0}^{z - x - 1} v^{k+1} \px[k]{x}[a] \px{x+k}[ai] \ax**{x + k + 1}[i] \]
where \( z \) is the pension age by law.


Individual Method: Consideration of the actual age of the entitled and their spouse 
(e.g., husband and wife).
Collective Method: Using general assumptions on probabilities to be married at death. Moreover,
assuming an age \( y(x) \) of the widow(er) at begin of the year of death of their deceased spouse.

Let 
\( M \): remaining lifetime of an x-year old man (in complete years)
\( N \): remaining lifetime of a y -year old woman (in complete years)

For the individual method, the settlement value is given by
\[ B = v^{M+1} \ax**{\angl{N - M}} \mathbf{1}_{\{N > M\}} \]
\[ \mathbb{E}[B] = \ax**{x}[w] = \sum_{m \geq 0} v^{m+1} \px[m]{x} \qx{x + m} \cdot \px[m+1]{y} \cdot \ax**[w]{y + m + 1} \]
A more involved expression for this can be given with
\[ \ax**{xy} = \mathbb{E}[\ax**{\angl{\min(M, N) + 1}}]\]
then 
\[ \ax**{x}[w] = \ax**{y} - \ax**{xy} \]

For the collective method, introduce
\( H \): Indicator variable, which states whether the man has a widow at death (then H = 1)
with associated probability \( h_x \) of a man aged \( x \) having a widow at death \( x \).
\( N \): remaining lifetime of a woman who is at the beginning of the year of death of her husband
\( y_{x+M} \)-years old.

The settlement value is
\[ B = v^{M+1} H \ax**{\angl N} \]
and 
\[ \mathbb{E}[B] = \sum_{m \geq 0} \sum_{n \geq 1} v^{m+1} \ax**{\angl n} \mathbb{P}[M = m, H = 1, N = n] \]
\[ = \ax**{x}[w] = \sum_{m \geq 0} v^{m+1} \px[m]{x} \qx{x+m} h_{x+m} \cdot \px[\frac{1}{2}]{y_{x+m} + \frac{1}{2}}[w] \cdot \ax**{y_{x+m}+1}[w] \]


Give expressions for obligations calculated with the Heubeck Richttafeln

Let \(k \in \mathbb{N}_{0}\) and \( \gamma = 0 \) be the event of receiving (legally required) pension due to old age, \\
\( \actsymb[k]{L}{x}[(0)] \): expected value of benefits which are obtained by reaching the year \(x + k\) in the main
population, discounted for the beginning of the year. \\
\( \actsymb[k]{L}{x}[(\gamma)] \) : expected value of benefits which are obtained by eliminated in the year \((x + k, x + k + 1] \) 
due to the cause \(\gamma\), discounted for the beginning of the year. \\
\(R_k\) : Amount of annuity in case of pension event after \(k\) years of service \\
\(w\) : Amount of dependent survivor’s pension as a percentage of the main annuity \\

Then 
\[
  \actsymb[k]{L}{x}[(0)] = \begin{cases}
    0, & k < z - x, \\
    R_k(\ax**{z}[r] + \ax**{z}[rw]) & k = z - x
  \end{cases}
\]
where \( \ax**{z}[r] \) is the exp. value of a lifelong annuity of 1 payable in advance for a pensioneer of required age \(z\).
and \(\ax**{z}[rw] \) the exp. value of an entitlement for a pensioned widow of age \(z\). 

\[
  \actsymb[k]{L}{x}[(1)] = \begin{cases}
    R_k v^{\frac{1}{2}}(\ax**{x + k \frac{1}{2}}[i] + w \ax**{x + k + \frac{1}{2}}[iw]), & k < z - x, \\
    0 & k = z - x
  \end{cases}
\]
\[
  \actsymb[k]{L}{x}[(2)] = \begin{cases}
    w R_k v^{\frac{1}{2}} h_{x+k} \ax**{y_{x + k} + \frac{1}{2}}[w], & k < z - x, \\
    0 & k = z - x
  \end{cases}
\]
for mixed obligations the exp. value of all benefits is
\[ \actsymb[k]{\hat{L}}{x} =  \actsymb[k]{L}{x}[(0)] + \sum_\gamma \qx{x+k}[(\gamma)]\actsymb[k]{L}{x}[(\gamma)] \]

And the settlement value is
\[ \actsymb[m]{B}{x} = \sum_{k\geq 0} v^k \px[k]{x+m} \cdot \actsymb[m + k]{\hat{L}}{x} \]


Cash value (annuity) of a lifelong annuity of amount \(1\), payable in \(t\) installments p.a. in advance, to an \(x\)-year old
pensioner
\[\ax**[][(t)]{x}[r] = \sum_{k\geq 0} v^k \px[k]{x}[r] \cdot \actsymb[m + k][(t)]{\hat{L}}{x} = \ax**{x}[r] - k^{(t)} \]
where 

\[\actsymb[m + k][(t)]{\hat{L}}{x} = \frac{1}{t} \sum_{m = 0}^{t-1} v^{\frac{m}{t}} \px[\frac{m}{t}]{x+k} =  1 - k^{(t)}(1 - v \px{x}[r]) \]
\(k^{(t)}\): Value at the end of the pension entry year for benefits payable in this first year of pension \\
This is independent of \(x\) but can still be evaluated from it, via  
\(k^{(t)} = \ax**{x} - \ax**{x}[(t)]\) or \(k^{(t)} = \frac{1 + i}{t} \sum_{m = 0}^{t-1} \frac{m}{t + m i} \) where \( i \) is the interest rate. \\

Opposed to annuities which are evaluated for currently paid benefits, entitlements are evaluated for paid benefits in the future and 
do not require adjustments for inter-annual payments according to the invariance principle.

Entitlement on Disability Pension
\[ \ax**{x}[ai] =  \ax**[][(t)]{x}[ai] = \sum_{k \geq 0} \px[k]{x}[a] \cdot \actsymb[][(t)]{\hat{L}}{x+k}[ai] \]
where 
\[
  \actsymb[][(t)]{L}{x+k}[ai] = \begin{cases}
    \actsymb{i}{x} v^{\frac{1}{2}} \ax**[][(t)]{x + \frac{1}{2}}[i], & x < z, \\
    0 & x = z
  \end{cases}
\]
\[
  \ax**[][(t)]{x + \frac{1}{2}}[i] = v^{\frac{1}{2}} \px[\frac{1}{2}]{x + \frac{1}{2}}[i] \ax**{x+1}[i] 
  = v^{\frac{1}{2}} \frac{1 - \qx{x}[i]}{1 - \frac{1}{2} \px{x}[i]} \ax**{x+1}[i] 
\]

As usual, for constant premiums paid on a recurrent basis (instead of a one time sum) we can determine them 
for a premium period of length \( n \) (paid in advance) by
\[ \actsymb[0]{B}{x} = \sum_{k=0}^{n-1} v^k \px[k]{x} \cdot P_x = \ax**{x:\angl n} P_x \]
so \( P_x = \frac{\actsymb[0]{B}{x}}{\ax**{x:\angl n}} \).

\end{document}
